% ============================================================
%  Openbook / Formelsammlung Template — strikt nach my_dhbw
%  Stil: 2-spaltig, kompakt, scriptsize Tabellen, dünne Regeln.
%  Renderer: lualatex (zweimal, wg. Unicode-Inhalt)
% ============================================================
\documentclass[10pt, a4paper]{article}

\usepackage[ngerman]{babel}
\usepackage{fontspec}
\setmainfont[
  Path=/usr/share/fonts/TTF/,
  Extension=.ttf,
  BoldFont=DejaVuSans-Bold,
  ItalicFont=DejaVuSans-Oblique,
  BoldItalicFont=DejaVuSans-BoldOblique
]{DejaVuSans}
\setsansfont[
  Path=/usr/share/fonts/TTF/,
  Extension=.ttf,
  BoldFont=DejaVuSans-Bold,
  ItalicFont=DejaVuSans-Oblique,
  BoldItalicFont=DejaVuSans-BoldOblique
]{DejaVuSans}
\setmonofont[Path=/usr/share/fonts/TTF/,Extension=.ttf]{DejaVuSansMono}

\usepackage[top=0.8cm, bottom=0.8cm, left=0.8cm, right=0.8cm, footskip=0.4cm]{geometry}
\usepackage{amsmath, amssymb}
\usepackage{multicol}
\usepackage{xcolor}
\usepackage{titlesec}
\usepackage{enumitem}
\usepackage{fancyhdr}
\usepackage{lastpage}
\usepackage{microtype}
\usepackage{etoolbox}

\usepackage{booktabs}
\usepackage{array}
\usepackage{tabularx}
% ltablex/longtable bewusst NICHT geladen — kollidiert mit multicols.
\usepackage{colortbl}
\usepackage[most]{tcolorbox}
\usepackage{listings}
\usepackage[normalem]{ulem}
\usepackage{adjustbox}
\usepackage[colorlinks=false]{hyperref}

\renewcommand{\familydefault}{\sfdefault}

% ── Theme (my_dhbw accent) ────────────────────────────────────
\definecolor{accent}{RGB}{0, 60, 113}
\definecolor{primaryblue}{RGB}{0, 60, 113}
\definecolor{codebg}{RGB}{245, 245, 245}
\definecolor{figurebg}{RGB}{232, 241, 255}
\definecolor{tablehintbg}{RGB}{240, 255, 240}
\definecolor{solutionbg}{RGB}{242, 255, 242}
\definecolor{rulegray}{RGB}{180, 180, 180}
\definecolor{tableheadbg}{RGB}{0, 60, 113}

% ── Header/Footer — wie my_dhbw formelsammlung.tex ────────────
\pagestyle{fancy}
\fancyhf{}
\renewcommand{\headrulewidth}{0pt}
\renewcommand{\footrulewidth}{0.3pt}
\fancyfoot[L]{\tiny\color{accent}\nichttitle}
\fancyfoot[R]{\tiny\color{accent}\thepage\,/\,\pageref{LastPage}}

% ── Typografie — wie formelsammlung.tex ───────────────────────
\titleformat{\section}
    {\color{accent}\footnotesize\bfseries}{}{0pt}{}
    [\vspace{-2pt}\color{accent!50}\rule{\linewidth}{0.4pt}]
\titleformat{\subsection}
    {\scriptsize\bfseries\color{accent!85!black}}{}{0pt}{}{}
\titleformat{\subsubsection}
    {\scriptsize\itshape\color{accent!60!black}}{}{0pt}{}{}
\titlespacing*{\section}{0pt}{4pt}{1pt}
\titlespacing*{\subsection}{0pt}{3pt}{0pt}
\titlespacing*{\subsubsection}{0pt}{2pt}{0pt}

\setlength{\parindent}{0pt}
\setlength{\parskip}{1pt}
\renewcommand{\arraystretch}{1.0}
\setlist{noitemsep, topsep=0pt, parsep=0pt, leftmargin=1.1em}

\setlength{\columnsep}{6mm}
\setlength{\columnseprule}{0.2pt}

% Globaler scriptsize-Body. Tabellen zusätzlich auf scriptsize zwingen.
\AtBeginEnvironment{tabularx}{\scriptsize}
\AtBeginEnvironment{tabular}{\scriptsize}

% ── Boxen (kompakt) ───────────────────────────────────────────
\newtcolorbox{figurebox}[1]{
  colback=figurebg, colframe=accent!50,
  title={\scriptsize Abbildung: #1}, fonttitle=\scriptsize\bfseries,
  fontupper=\scriptsize, sharp corners,
  boxsep=2pt, left=3pt, right=3pt, top=2pt, bottom=2pt,
  before skip=3pt, after skip=3pt
}
\newtcolorbox{tablehintbox}[1]{
  colback=tablehintbg, colframe=green!50!black!50,
  title={\scriptsize Tabelle: #1}, fonttitle=\scriptsize\bfseries,
  fontupper=\scriptsize, sharp corners,
  boxsep=2pt, left=3pt, right=3pt, top=2pt, bottom=2pt,
  before skip=3pt, after skip=3pt
}
\newtcolorbox{solutionbox}[1]{
  colback=solutionbg, colframe=green!60!black!60,
  title={\scriptsize #1}, fonttitle=\scriptsize\bfseries,
  fontupper=\scriptsize, sharp corners,
  boxsep=2pt, left=3pt, right=3pt, top=2pt, bottom=2pt,
  before skip=3pt, after skip=3pt
}

\lstset{
  basicstyle=\ttfamily\tiny,
  backgroundcolor=\color{codebg},
  breaklines=true,
  frame=none,
  xleftmargin=2pt, xrightmargin=2pt,
  aboveskip=2pt, belowskip=2pt,
  keepspaces=true
}

\providecommand{\nichttitle}{Formelsammlung}
\providecommand{\nichtsubtitle}{}

\renewcommand{\nichttitle}{Betriebssysteme}
\renewcommand{\nichtsubtitle}{Formelsammlung}


\begin{document}

\begin{center}
{\small\bfseries\color{accent}\nichttitle}
\end{center}
\vspace{-6pt}

\scriptsize
\begin{multicols}{2}

% ── BODY ──────────────────────────────────────────────────────

\section{Grundlagen}

\textbf{BS}: Schnittstelle HW↔SW; steuert Programmabwicklung, regelt Rechnerbetrieb, stellt Dienste bereit. Kern = \textbf{Kernel}.


\textbf{Aufgaben}: Prozess-, Speicher-, Geräte-, Datei-, Energieverwaltung.



{\small
\begin{adjustbox}{max width=\linewidth}
{\scriptsize
\begin{tabular}{ll}
\toprule
\textbf{Programmart} & \textbf{Zweck} \\
\midrule
Anwendung & löst Benutzer-Probleme \\
Systemprogramm & verwaltet Rechnerbetrieb \\
BS & elementare Dienste \\
\bottomrule
\end{tabular}
}
\end{adjustbox}
}


\textbf{Schnittstellen}: GUI (grafisch), CLI (Kommandozeile), TUI (zeichen), VUI (Sprache), NUI (natürlich).



\section{Rechnerarchitektur}

\textbf{Hierarchie} (schnell↓/groß↑): Register → Cache → RAM/ROM → Sekundär.


\textbf{CPU}: Steuerwerk (Befehl), ALU/Rechenwerk (Daten + Flags + Akku); Register IR/MBR/PC/MAR; Systembus = Adress+Daten+Steuer.


\textbf{Registerbreite}: 4/8/16/32/64 Bit; 64-Bit-OS adressiert 2⁶⁴.



\section{Modi \& Syscall}


{\small
\begin{adjustbox}{max width=\linewidth}
{\scriptsize
\begin{tabular}{lll}
\toprule
\textbf{} & \textbf{User} & \textbf{Kernel} \\
\midrule
Privileg & nicht-priv. & privilegiert \\
Befehlssatz & eingeschränkt & beliebig \\
Speicher & kein Kernel & gesamt + alle BM \\
\bottomrule
\end{tabular}
}
\end{adjustbox}
}


\begin{itemize}
  \item Kernel→User: unproblematisch via BS
  \item User→Kernel: nur über \textbf{Syscall}
\end{itemize}

\textbf{Syscall-Ablauf}: Aufruf → Kernel-Mode → Sicherheitsprüfung → Ausführen/Ablehnen → Rückkehr User-Mode.


Bsp.: \texttt{open}, \texttt{read}, \texttt{write}, \texttt{close}.



\section{Betriebsmittel}


{\small
\begin{adjustbox}{max width=\linewidth}
{\scriptsize
\begin{tabular}{lll}
\toprule
\textbf{Art} & \textbf{Eigenschaft} & \textbf{Bsp.} \\
\midrule
Entziehbar & jederzeit ohne Schaden & CPU \\
Nicht entziehbar & bis Beendigung & Drucker \\
Exklusiv & max. 1 Prozess & — \\
Gemeinsam & quasi-gleichzeitig & Festplatte \\
\bottomrule
\end{tabular}
}
\end{adjustbox}
}


\textbf{BS-Kategorien}: Mainframe (z/OS), Server, PC, Echtzeit, Embedded, Tablet, Handheld, Smartcard.



\section{Prozess}

\textbf{Prozess}: Programm zur Laufzeit; dynamische Aktionsfolge; Adressraum = Userspace + Kernelspace.


\textbf{Programm}: statische Verfahrensvorschrift.


\textbf{Erzeugung}: Code laden → PID zuordnen → PCB anlegen. Syscalls: Unix \texttt{fork}, Windows \texttt{CreateProcess}.


\textbf{Kenndaten}: PID (≥0, eindeutig), PPID, UID/GID, EUID/EGID, PRI, Zustand, CPU-Zeit.



{\small
\begin{adjustbox}{max width=\linewidth}
{\scriptsize
\begin{tabular}{lll}
\toprule
\textbf{BS} & \textbf{PID 0} & \textbf{PID 1} \\
\midrule
Unix & Idle & Init \\
Windows & nicht vergeben & nicht vergeben \\
\bottomrule
\end{tabular}
}
\end{adjustbox}
}



\subsection{PCB-Inhalt}
Zustand, PC, PID, init+dyn. Priorität, CPU-Zeit, BM, Register, HW-Kontext (PC/IR/Stack/Flags, MMU, Stack/Heap/Code).


\textbf{Kontextwechsel}: HW-Kontext A sichern (PCB\_A) → HW-Kontext B laden (PCB\_B).



\subsection{Zustände}
\textbf{Bereit, Aktiv, Blockiert, Beendet}.

\begin{enumerate}
  \item Aktivieren (Bereit→Aktiv)
  \item Deaktivieren (Aktiv→Bereit)
  \item Blockieren (Aktiv→Blockiert) — wartet I/O
  \item Deblockieren (Blockiert→Bereit)
  \item Terminieren
\end{enumerate}

\textbf{Sonderzustände}:

\begin{itemize}
  \item \textbf{Zombie}: beendet, noch in Tabelle; Eltern muss austragen
  \item \textbf{Daemon}: Hintergrund, terminalunabhängig
\end{itemize}


\subsection{Thread}

\textbf{Thread}: kleinste Ausführungseinheit; Prozess = ≥1 Thread (Main); Threads teilen Code/Daten/Heap.



{\small
\begin{adjustbox}{max width=\linewidth}
{\scriptsize
\begin{tabular}{lll}
\toprule
\textbf{} & \textbf{Prozess} & \textbf{Thread} \\
\midrule
Gewicht & schwer & leicht \\
Adressraum & eigen & gemeinsam \\
Speicherschutz & ja & nein \\
Kontextwechsel & aufwendig & schnell \\
\bottomrule
\end{tabular}
}
\end{adjustbox}
}


\textbf{Vorteile}: schneller KW, Modularisierung. \textbf{Nachteile}: Sync nötig, kein Schutz, schwerer zu programmieren.


\textbf{Zustände}: Laufend, Bereit, Wartend.



{\small
\begin{adjustbox}{max width=\linewidth}
{\scriptsize
\begin{tabular}{ll}
\toprule
\textbf{System} & \textbf{Prio-Bereich} \\
\midrule
Windows & 0 System; 1–15 Timesharing; 16–31 Realtime \\
Linux & 0–99 Realtime; 100–139 Timesharing \\
\bottomrule
\end{tabular}
}
\end{adjustbox}
}


\textbf{GC}: niedrige Prio, läuft bei Leerlauf/Knappheit.



\section{Scheduling}

\textbf{Scheduler}: wer bekommt CPU. \textbf{Dispatcher}: führt KW aus. \textbf{Prozessmanager} = Scheduler + Dispatcher.


\textbf{Horizonte}: kurzfristig (CPU), mittelfristig (Speicher), langfristig (eingehend).



{\small
\begin{adjustbox}{max width=\linewidth}
{\scriptsize
\begin{tabular}{ll}
\toprule
\textbf{Klasse} & \textbf{Bedeutung} \\
\midrule
Non-preemptive & nicht-verdrängend; nur bei Block/Freigabe/Ende \\
Preemptive & verdrängend; jederzeit entziehbar \\
\bottomrule
\end{tabular}
}
\end{adjustbox}
}


\textbf{Ziele}: Fairness, Effizienz, min. Antwort-/Warte-/Durchlaufzeit, max. Durchsatz, Vorhersehbarkeit.



{\small
\begin{adjustbox}{max width=\linewidth}
{\scriptsize
\begin{tabular}{ll}
\toprule
\textbf{Begriff} & \textbf{Definition} \\
\midrule
Antwortzeit & Wartezeit Anwender auf Bearbeitung \\
Wartezeit & Σ aller Wartezeiträume \\
Durchsatz & Prozesse/Zeit \\
Durchlaufzeit & Bedienzeit + Warten \\
Bedienzeit & CPU-Haltezeit \\
\bottomrule
\end{tabular}
}
\end{adjustbox}
}



\subsection{Algorithmen}


{\small
\begin{adjustbox}{max width=\linewidth}
{\scriptsize
\begin{tabular}{lll}
\toprule
\textbf{Verfahren} & \textbf{Typ} & \textbf{Prinzip} \\
\midrule
FCFS & Batch & Reihenfolge Eintreffen \\
SJF/SPF/SPN & Batch & kürzester zuerst \\
SRTN & Batch, preempt & preemptive SJF \\
Round Robin & Dialog & FCFS + Quantum \\
PS & Dialog & höchste Prio zuerst \\
FSS & Dialog & 1/n CPU; schlechtestes Verhältnis zuerst \\
Lottery & Dialog & n Verlosungen/s, m ms CPU \\
EDF/MDF & RT & kleinste Deadline \\
Polled Loop & RT & zyklische Abfrage \\
Interrupt-gest. & RT & ISR auf Ereignis \\
\bottomrule
\end{tabular}
}
\end{adjustbox}
}


\begin{quote}
SJF/SRTN: Bedienzeit unbekannt; \textbf{Starvation} möglich.
\end{quote}

\textbf{Quantum}: typ. \textbf{10–200 ms}; zu kurz=Overhead, zu lang=Verzögerung; höhere Taktrate→kürzer; E/A-intensiv höher, rechenintensiv kürzer.


\textbf{Priorität}: statisch + dynamisch (kombiniert).



{\small
\begin{adjustbox}{max width=\linewidth}
{\scriptsize
\begin{tabular}{ll}
\toprule
\textbf{RT-Typ} & \textbf{Bedeutung} \\
\midrule
Hard & feste Schranke; Fehler bei Verspätung \\
Soft & weiche Schranke; Verspätung lästig \\
\bottomrule
\end{tabular}
}
\end{adjustbox}
}



\subsection{Multi-Level}


{\small
\begin{adjustbox}{max width=\linewidth}
{\scriptsize
\begin{tabular}{ll}
\toprule
\textbf{Prio} & \textbf{Klasse} \\
\midrule
0 & System \\
1 & Interaktiv \\
2 & Allgemein \\
4 & Batch \\
\bottomrule
\end{tabular}
}
\end{adjustbox}
}


Pro Stufe Queue; innerhalb RR. \textbf{Multi-Level-Feedback}: Klassenwechsel möglich (Praxis).



\section{Interrupts}

\textbf{Interrupt}: Unterbrechung sequenzieller Ausführung; Kontrolle an \textbf{ISR}.


\textbf{Befehlszyklus}: Fetch → Execute → IRQ-Prüfung → ggf. PC=ISR-Adresse.



{\small
\begin{adjustbox}{max width=\linewidth}
{\scriptsize
\begin{tabular}{lll}
\toprule
\textbf{Verfahren} & \textbf{+} & \textbf{−} \\
\midrule
Polling & leicht & CPU-Verschwendung \\
Interrupt & effizient & komplex \\
\bottomrule
\end{tabular}
}
\end{adjustbox}
}



{\small
\begin{adjustbox}{max width=\linewidth}
{\scriptsize
\begin{tabular}{lll}
\toprule
\textbf{Klasse} & \textbf{Eigenschaft} & \textbf{Bsp.} \\
\midrule
Asynchron & unvorhersehbar & HW-IRQ (Platte, Netz) \\
Synchron (Exception) & reproduzierbar & Syscall, Trap, Fault \\
\bottomrule
\end{tabular}
}
\end{adjustbox}
}


\begin{itemize}
  \item \textbf{Fault}: Unterbrechung \textbf{vor} Ausführung (Page Fault)
  \item \textbf{Trap}: Unterbrechung \textbf{nach} Ausführung (Div/0)
\end{itemize}


\subsection{Komponenten}
\begin{itemize}
  \item \textbf{IC} (Interrupt Controller): Signale → CPU
  \item \textbf{ISR}: pro Interrupt-Typ; bei Geräten im Treiber
  \item \textbf{IVT}: Kernelspeicher; Eintrag = ISR-Adresse, Index = Interrupt-Nr.
  \item \textbf{IL}: Priorität pro Interrupt
\end{itemize}

\textbf{Ablauf}: IRQ (wenn IL erlaubt) → Status auf Stack → ISR via IVT (Kernel) → Bestätigung IC → Status wiederherstellen.


\begin{quote}
\textbf{Merke}: Ergebnis darf nicht von IRQ-Auftreten abhängen. Gleich-/niedriger-priorisierte werden während Bearbeitung nicht angenommen.
\end{quote}


\subsection{Maskierung}
\begin{itemize}
  \item \textbf{IMR}: Bit=1 → Interrupt aus; IC wiederholt bis angenommen
\end{itemize}


{\small
\begin{adjustbox}{max width=\linewidth}
{\scriptsize
\begin{tabular}{ll}
\toprule
\textbf{Typ} & \textbf{Bedeutung} \\
\midrule
MI & maskierbar \\
NMI & nicht maskierbar \\
Präzise & komplett oder gar nicht \\
Unpräzise & teilweise möglich \\
\bottomrule
\end{tabular}
}
\end{adjustbox}
}



\section{Synchronisation}

\textbf{Nebenläufigkeit}: quasi-parallel (1 Kern) / echt-parallel (mehrere).


\textbf{Probleme}: Blockieren, Verhungern, Verklemmung.


\textbf{Race Condition}: Ergebnis hängt von Zugriffsreihenfolge ab. Voraussetzung: Nebenläufigkeit + gemeinsame BM.


\textbf{Lost Update}: A=0, B=0, A schreibt, B schreibt → A verloren.



\subsection{Kritischer Abschnitt}

\textbf{KA}: nur 1 Prozess gleichzeitig.


\textbf{Atomar}: logisch nicht unterbrechbar.


\textbf{Mutex-Bedingungen}:

\begin{enumerate}
  \item Mutual Exclusion
  \item Keine Annahmen über Geschwindigkeit/Anzahl CPUs
  \item Außerhalb darf niemand andere blockieren
  \item Fairness — jeder Wartende kommt rein
\end{enumerate}

\textbf{Implementierung}: Busy Waiting (CPU-Verschwendung) | Schlafen+Aufwecken.



\subsection{Konzepte}
\begin{enumerate}
  \item Software (selten)
  \item Interrupts sperren — bei MP unpraktisch
  \item HW-Befehle (TSL/TAS)
\end{enumerate}

\textbf{Sperrvariable}: 0=frei, 1=belegt. \textbf{Spinlock}: Warteschleife (kurze Wartezeit).


\textbf{TSL/TAS}: Test+Set ununterbrechbar in 1 Maschinenbefehl, 1 Speicherzyklus, CPU blockiert Bus, MP-tauglich.



\subsection{Semaphor}

Zähler + Warteschlange.



{\small
\begin{adjustbox}{max width=\linewidth}
{\scriptsize
\begin{tabular}{ll}
\toprule
\textbf{Op} & \textbf{Wirkung} \\
\midrule
\textbf{P() / Down} & Z>0: Z−−; sonst suspendieren \\
\textbf{V() / Up} & Z++; ggf. aufwecken \\
\bottomrule
\end{tabular}
}
\end{adjustbox}
}


\begin{itemize}
  \item \textbf{Mutex / Binär}: Z∈\{0,1\}
  \item \textbf{Zählsemaphor}: max. n im KA
\end{itemize}

\texttt{P();KA;V();} | \texttt{V();...P();} → alle drin | \texttt{P();...P();} → niemand


\begin{quote}
Semaphore fehleranfällig.
\end{quote}


\subsection{Monitor}

Objekt mit gem. Daten + Methoden + Türsteher (1 aktiv); Compiler erzeugt Sync.


\textbf{Bedingungsvariablen} mit \texttt{wait(cond)} (warten, Monitor verlassen) / \texttt{signal(cond)} (aufwecken).



{\small
\begin{adjustbox}{max width=\linewidth}
{\scriptsize
\begin{tabular}{lll}
\toprule
\textbf{Konzept} & \textbf{Typ} & \textbf{Verhalten} \\
\midrule
Signal-and-Continue & Mesa & Signalisierender arbeitet weiter \\
Signal-and-Wait & Hoare & Signalisierender verlässt sofort \\
\bottomrule
\end{tabular}
}
\end{adjustbox}
}



\section{Deadlock}

\textbf{Deadlock}: Prozesse halten BM und warten auf weitere → Zyklus.



{\small
\begin{adjustbox}{max width=\linewidth}
{\scriptsize
\begin{tabular}{ll}
\toprule
\textbf{Begriff} & \textbf{Definition} \\
\midrule
Blockieren & P2 wartet auf von P1 gehaltene Ressource \\
Verhungern & trotz Bereitschaft keine CPU \\
\bottomrule
\end{tabular}
}
\end{adjustbox}
}



\subsection{4 Bedingungen (alle gleichzeitig)}
\begin{enumerate}
  \item Mutual Exclusion
  \item Hold and Wait
  \item No Preemption
  \item Circular Wait
\end{enumerate}


\subsection{Strategien}
\begin{enumerate}
  \item \textbf{Erkennen}: Betriebsmittelgraph + Zyklussuche
  \item \textbf{Ignorieren} (Vogel-Strauß): extern abbrechen
  \item \textbf{Vermeiden}: BS so bauen, dass unmöglich
  \item \textbf{Verhindern}: Anforderung nur ohne zukünftigen DL
\end{enumerate}


\subsection{Klassische Probleme}
\begin{itemize}
  \item \textbf{Producer-Consumer} (Bounded Buffer)
  \item \textbf{Reader-Writer}
  \item \textbf{Dining Philosophers}: 5 Phil., 5 Gabeln, 2 nötig, max 2 essen, niemand verhungert
\end{itemize}


\section{Speicher — Grundlagen}

\textbf{Realer AR}: physisch verbaut. \textbf{Phys. Adresse}: im phys. Speicher. \textbf{AR-Größe}: z.B. 4 GB=\{0…2³²−1\}.


\textbf{Inhalte}: statisch (Code, Konstanten, init/uninit Daten) | dynamisch (Heap↑, Stack↓).


\textbf{Techniken}: Mono → Multi feste Partitionen → Multi variable + Swapping → Multi virtuell + Paging.



\section{Virtueller Speicher}

\textbf{Virt. AR}: pro Prozess; aktuell genutzt im RAM, Rest auf Sekundärspeicher; Bedarf > RAM erlaubt.


\textbf{Vorteile}: Programme>RAM, mehr Programme im RAM, AR>physisch, externe Fragmentierung untergeordnet.


\textbf{Page}: virt. Block. \textbf{Page Frame}: RAM-Block. Gleich groß: typ. \textbf{1, 4, 8, 16, 64 KiB}.


\textbf{Paging Area / Schattenspeicher}: HDD-Bereich für ausgelagerte Pages.



{\small
\begin{adjustbox}{max width=\linewidth}
{\scriptsize
\begin{tabular}{lll}
\toprule
\textbf{} & \textbf{Swapping} & \textbf{Paging} \\
\midrule
Granularität & ganzer Prozess & Teile \\
MMU & ohne & mit \\
\bottomrule
\end{tabular}
}
\end{adjustbox}
}



\section{MMU}

CPU-Funktionseinheit für virt→real; ersetzt Basis/Limit; Pfad: \texttt{CPU→virt→MMU→real→RAM}; greift auf Seitentabellen zu.



\section{Seitentabellen}

Pro virt. AR eine; enthält \textbf{Verweise auf Frames}, nicht Daten.


\textbf{Adressaufbau}: virt = \texttt{i | d} → real = \texttt{f | d}.



{\small
\begin{adjustbox}{max width=\linewidth}
{\scriptsize
\begin{tabular}{ll}
\toprule
\textbf{Feld} & \textbf{Bedeutung} \\
\midrule
P/A & Present/Absent \\
Protection (C/R/M) & r/w/x \\
Frame f & Verweis \\
M (Dirty) & verändert \\
R (Reference) & Zugriff \\
C & Caching \\
\bottomrule
\end{tabular}
}
\end{adjustbox}
}



{\small
\begin{adjustbox}{max width=\linewidth}
{\scriptsize
\begin{tabular}{llll}
\toprule
\textbf{Art} & \textbf{Aufbau} & \textbf{+} & \textbf{−} \\
\midrule
Einstufig & Index→Frame & direkt & enormer Speicher (32-Bit/4KiB→>10⁶) \\
Mehrstufig & `i1\textbackslash{} & i2\textbackslash{} & d` (32-Bit: 10+10+12) \\
Invertiert & pro Frame; `pid\textbackslash{} & i\textbackslash{} & f`+Kollisionsliste; 12-Bit Hash \\
\bottomrule
\end{tabular}
}
\end{adjustbox}
}


\textbf{TLB}: CPU-Cache aktiver (virt→frame). Hit→direkt; Miss→Tabellenlookup.



\section{Page Fault \& Ersetzung}

\textbf{Page Fault}: Trap, Seite nicht im Frame; Bearbeitung im Kernel.



{\small
\begin{adjustbox}{max width=\linewidth}
{\scriptsize
\begin{tabular}{ll}
\toprule
\textbf{Begriff} & \textbf{Definition} \\
\midrule
Eingelagert & in Frame \\
Ausgelagert & nicht in Frame \\
Modifiziert & M=1 \\
\bottomrule
\end{tabular}
}
\end{adjustbox}
}


\begin{quote}
Modifizierte Seiten nicht ohne Sicherung ersetzen.
\end{quote}


{\small
\begin{adjustbox}{max width=\linewidth}
{\scriptsize
\begin{tabular}{lll}
\toprule
\textbf{} & \textbf{Demand} & \textbf{Prepaging} \\
\midrule
Politik & nur auf Anforderung & spekulativ \\
+ & nur Benötigtes & nutzt Lokalität \\
\bottomrule
\end{tabular}
}
\end{adjustbox}
}


\textbf{Ersetzungsablauf}: PF → kein freier Frame → SEA → Opfer A → wenn modifiziert: in Paging Area → B in Frame X.



\section{SEA (Seitenersetzungsalgorithmen)}


{\small
\begin{adjustbox}{max width=\linewidth}
{\scriptsize
\begin{tabular}{ll}
\toprule
\textbf{Algo} & \textbf{Prinzip} \\
\midrule
\textbf{Optimal (Belady)} & längste Zukunft ungenutzt — nur Referenz \\
\textbf{FIFO} & älteste eingelagert; R nicht nötig; ignoriert Nutzung \\
\textbf{NRU} & 4 Klassen aus (R,M); R period. reset (\textasciitilde{}20 ms); Klasse 1 zuerst \\
\textbf{Second Chance} & FIFO+R: R=1→R:=0+ans Ende; R=0→raus \\
\textbf{Clock Page} & Second Chance zirkulär mit Zeiger; max 1 Runde \\
\textbf{LRU} & längste ungenutzte; Matrix (n×n) oder Zeitstempel; teuer \\
\textbf{NFU} & Zähler kleinster raus; Z+=R, R:=0; alte Zähler problematisch \\
\textbf{Aging} & NFU + Bit-Shift vor Add; R = MSB; ≈ LRU \\
\textbf{Working Set} & aktuell nötige Seiten ständig im RAM; Δ = Distanz \\
\bottomrule
\end{tabular}
}
\end{adjustbox}
}


\textbf{NRU-Klassen}:


{\small
\begin{adjustbox}{max width=\linewidth}
{\scriptsize
\begin{tabular}{llll}
\toprule
\textbf{Klasse} & \textbf{R} & \textbf{M} & \textbf{Verhalten} \\
\midrule
1 & 0 & 0 & keine Sicherung \\
2 & 0 & 1 & schreiben \\
3 & 1 & 0 & wahrsch. wieder \\
4 & 1 & 1 & schreiben+wieder \\
\bottomrule
\end{tabular}
}
\end{adjustbox}
}


\textbf{LRU-Matrix}: Zugriff k → Zeile k=1, Spalte k=0; Max=zuletzt; Min=Opfer.


\textbf{Pseudo-LRU}: Clock, Second Chance.


\textbf{Lokalität (80/20)}: kleiner Code-Teil aktiv.



\section{Lokalität}


{\small
\begin{adjustbox}{max width=\linewidth}
{\scriptsize
\begin{tabular}{lll}
\toprule
\textbf{} & \textbf{Zeitlich} & \textbf{Räumlich} \\
\midrule
Bezug & Zeit & Abstand \\
Aussage & Adresse bald wieder & Nachbar bald \\
Optim. & im Cache halten & Nachbarn mitladen \\
\bottomrule
\end{tabular}
}
\end{adjustbox}
}



\section{Fragmentierung}


{\small
\begin{adjustbox}{max width=\linewidth}
{\scriptsize
\begin{tabular}{lll}
\toprule
\textbf{Art} & \textbf{Wo} & \textbf{Bedeutung} \\
\midrule
Intern & innerhalb Block & Block>Bedarf (Ø 50\% letzte Seite) \\
Extern & zwischen Blöcken & Lücken einzeln zu klein \\
\bottomrule
\end{tabular}
}
\end{adjustbox}
}


\begin{itemize}
  \item RAM/HDD: beide; SSD: v.a. intern
  \item \textbf{Defrag}: Blöcke zusammenhängend → Performance ↑ (HDD)
\end{itemize}


\section{Placement}

\textbf{Verwaltung freier Bereiche}:

\begin{itemize}
  \item \textbf{Bitmap}: 1 Bit/Frame (1=belegt, 0=frei); freie = 0-Folgen
  \item \textbf{Verkettete Liste}
\end{itemize}


{\small
\begin{adjustbox}{max width=\linewidth}
{\scriptsize
\begin{tabular}{ll}
\toprule
\textbf{Strategie} & \textbf{Vorgehen} \\
\midrule
First-Fit & erster passend \\
Next-Fit & First ab letzter Pos \\
Best-Fit & kleinste passende \\
Worst-Fit & größte Lücke \\
Buddy & 2ᵏ-Halbierung \\
Quick-Fit & Buddy, nicht doppelt \\
\bottomrule
\end{tabular}
}
\end{adjustbox}
}



\subsection{Buddy}
\begin{itemize}
  \item 2ᵏ-Blöcke; Anforderung auf 2ᵏ aufgerundet
  \item Halbierung; benachbarte gleich-große freie = Buddies → Verschmelzen
  \item + gegen externe Frag; − interne Frag durch Aufrundung
\end{itemize}


\subsection{Quick-Fit}
Buddy ohne Verdopplung; Verschmelzen aufwändiger.



\section{Festplatten / Cluster}

\textbf{Cluster}: benachbarte Sektoren; kleinste Adressierungseinheit; ≥1 pro Datei.



{\small
\begin{adjustbox}{max width=\linewidth}
{\scriptsize
\begin{tabular}{ll}
\toprule
\textbf{Slack} & \textbf{Bedeutung} \\
\midrule
File Slack & ungenutzt im letzten Cluster \\
RAM Slack & EOF→Sektorende, mit RAM-Inhalt \\
Drive Slack & Sektoren bis Clustergrenze, mit HDD-Daten \\
\bottomrule
\end{tabular}
}
\end{adjustbox}
}


\begin{itemize}
  \item Linux: Zufallszahlen statt RAM/Drive
  \item Forensik: wertvoll
  \item Große Partition→große Cluster→Slack↑
\end{itemize}


\section{Geräteverwaltung}

Schnittstelle Geräte↔BS; Prozesse nur via Syscall.


\textbf{Aufgaben}: Kommandos senden, IRQ-Bearbeitung, Fehlerbehandlung.


\textbf{Architektur}: Northbridge (RAM, AGP) | Southbridge (PCI, S-ATA, SCSI, USB).


\textbf{Transferraten}:


{\small
\begin{adjustbox}{max width=\linewidth}
{\scriptsize
\begin{tabular}{ll}
\toprule
\textbf{Gerät} & \textbf{Rate} \\
\midrule
Keyboard & \textasciitilde{}10 B/s \\
Maus & \textasciitilde{}100 B/s \\
Laserdrucker & \textasciitilde{}100 KB/s \\
ISA & 16 MB/s \\
USB1/2/3 & 12 / 480 Mbit/s / 5 Gbit/s \\
Fast/Gigabit Ethernet & 100 Mbit/s / bis 10 Gbit/s \\
PCIe & 32 Gbit/s \\
ATA / SATA & 130 MB/s / \textasciitilde{}2 GB/s \\
\bottomrule
\end{tabular}
}
\end{adjustbox}
}



\subsection{Treiber}
\textbf{Aufgaben}: bekanntmachen, Controller init, Befehle übersetzen, Pufferung, ISR, Sync.



{\small
\begin{adjustbox}{max width=\linewidth}
{\scriptsize
\begin{tabular}{llll}
\toprule
\textbf{Klasse} & \textbf{Übertragung} & \textbf{Bsp.} & \textbf{Funktionen} \\
\midrule
Block & komplette Blöcke 512–32768 B, adressierbar & HDD, Band, CD/DVD & initDevice, readBlock, writeBlock, handleInterrupt \\
Zeichen & Datenstrom & Maus, Tastatur, NIC, ser/par/USB & initDevice, readChar, writeChar, handleInterrupt \\
\bottomrule
\end{tabular}
}
\end{adjustbox}
}



\subsection{Controller}
Register: Schreiben→Befehl; Lesen→Status. Syscalls: open/close/read/write.



{\small
\begin{adjustbox}{max width=\linewidth}
{\scriptsize
\begin{tabular}{llll}
\toprule
\textbf{Methode} & \textbf{Adressraum} & \textbf{+} & \textbf{−} \\
\midrule
\textbf{I/O-Port} & getrennt (8/16-Bit Portnr.) & einfach & spez. Port-Befehle \\
\textbf{Memory-Mapped I/O} & gemeinsam (oben) & Zugriff wie RAM, leichter Schutz & echt/gemappt aufwändig \\
\bottomrule
\end{tabular}
}
\end{adjustbox}
}


\begin{quote}
Beide kosten viel CPU-Zeit.
\end{quote}


\subsection{DMA}

Block-Transfer Controller↔RAM via DMA-Controller; CPU nur Anfang/Ende.


\textbf{+}: CPU-Entlastung, Geschwindigkeit↑, weniger IRQ (1 statt 1/Wort), spezialisiert.


\textbf{Ablauf (Lesen)}: CPU programmiert DMA (Quelle, Ziel, n) → Lese-Cmd an Platten-Ctrl → DMA fordert an → Platte→DMA→RAM → Abschluss → \textbf{1 IRQ} → ISR.



\section{Dateiverwaltung}

\textbf{CRUD}: Create, Read, Update, Delete.


\textbf{Dienste}: Datei/Katalog Erzeugen/Löschen/Kopieren, Verknüpfungen, Umbenennen.



\section{Dateisystem}

\textbf{Aufgaben}: CRUD, Frei/Belegt-Verwaltung, Strukturierung, Verzeichnisse, Attribute, Rechte.



{\small
\begin{adjustbox}{max width=\linewidth}
{\scriptsize
\begin{tabular}{ll}
\toprule
\textbf{Speicherung} & \textbf{Eigenschaft} \\
\midrule
n aufeinanderfolgende Bytes & Größenänderung→Verschiebung \\
Aufteilung in Blöcke & Standard, flexibel \\
\bottomrule
\end{tabular}
}
\end{adjustbox}
}



{\small
\begin{adjustbox}{max width=\linewidth}
{\scriptsize
\begin{tabular}{ll}
\toprule
\textbf{BS} & \textbf{Dateisysteme} \\
\midrule
Apple & DOS, ProDOS, MFS, HFS/HFSX/HFS+, APFS \\
Linux & ext, ext2, ext3, ext4, ReiserFS \\
Windows & FAT12/16/32, exFAT, VFAT, NTFS, ReFS \\
\bottomrule
\end{tabular}
}
\end{adjustbox}
}



\section{Dateisystemcheck}

\textbf{Optionen}: 1) immer konsistent halten | 2) Journal/Log.



\subsection{Journaling}

\begin{quote}
Log-Eintrag \textbf{immer vor} Änderung.
\end{quote}

\textbf{Transaktion}: atomar (komplett/nicht). Bsp.: Erzeugen/Löschen/Erweitern/Verkürzen/Umbenennen, Attribute.


\textbf{Inkonsistenzen}: Stromausfall, HW-Ausfall, Absturz (z.B. Verzeichniseintrag fehlt; Block belegt aber nicht markiert).


\textbf{Wiederherstellung}: beim Boot Protokoll abgleichen.



{\small
\begin{adjustbox}{max width=\linewidth}
{\scriptsize
\begin{tabular}{lll}
\toprule
\textbf{Art} & \textbf{Inhalt} & \textbf{Eigenschaft} \\
\midrule
Meta-Daten & Operationen ohne Inhalt & Datenverlust möglich \\
Vollständig/Data & inkl. Daten & doppeltes Schreiben, Original gesichert \\
\bottomrule
\end{tabular}
}
\end{adjustbox}
}


Meta-Daten: Verzeichnisreferenz, Rechte, Eigentümer, Blöcke.



\subsection{Log-structured}

Ganze Platte = Log; \textbf{Copy-on-Write}; Verweise umgebogen.


\textbf{Aufbau}: Log Head → Daten → Empty → Daten → Log Tail.


\textbf{Schreiben}: RAM-Puffer → periodisch als Segment ans Ende.


\textbf{Cleaner}: Zusammenfassung erstes Segment lesen → I-Node-Map prüfen → aktuell: in Speicher+neu schreiben | nicht: löschen+freigeben.


\textbf{Nachteile}: Springen Log↔Daten; doppeltes Schreiben (Data-J.); starke Fragmentierung (Log-struct.).



\section{Backup}


{\small
\begin{adjustbox}{max width=\linewidth}
{\scriptsize
\begin{tabular}{llll}
\toprule
\textbf{Art} & \textbf{Basis} & \textbf{Speicher} & \textbf{Restore} \\
\midrule
Voll & — & sehr hoch & nur Voll \\
Inkrementell & letztes Backup & niedrig & Voll + alle Inkremente \\
Differentiell & letztes Vollbackup & mittel & Voll + gewünschtes Diff \\
\bottomrule
\end{tabular}
}
\end{adjustbox}
}


\begin{itemize}
  \item Inkrementell: fehlt eines → Dateien fehlen
  \item Differentiell: einmal verändert → bei jedem Diff erneut
\end{itemize}


\section{Meltdown \& Spectre}

\textbf{Meltdown}: liest geschützten RAM via CPU-Designfehler; Google PZ 28.7.2017, public 3.1.2018; alle Intel ab 1995; \textbf{Seitenkanal über Cache}.



\subsection{Grundlagen}


{\small
\begin{adjustbox}{max width=\linewidth}
{\scriptsize
\begin{tabular}{ll}
\toprule
\textbf{Modus} & \textbf{Rechte} \\
\midrule
Kernel & uneingeschränkt, gesamte HW+Speicher \\
User & nur via API \\
\bottomrule
\end{tabular}
}
\end{adjustbox}
}


\textbf{Speicherisolation}: pro Prozess virt. AR + Page-Table; KW=Flush+Laden (teuer).


\textbf{Kernel-Einbettung}: Performance → Kernel-AR + ges. phys. Speicher in \textbf{jeder} User-PT, geschützt durch \textbf{Supervisor-Bit}. User→Kernel-Adresse = \textbf{Seg-Fault}.



\subsection{CPU-Optimierungen}

\textbf{Pipelining (5)}: FETCH | DECODE | LOAD | EXECUTE | WRITE.


\textbf{Out-of-Order}: Befehle umsortiert, wenn Daten warten.


\textbf{Speculative}: rät vorab; Fehl → verworfen, \textbf{aber Cache/TLB nicht zurückgesetzt}; Schutzregel gilt nicht → Exploit.



\subsection{Angriff}

\textbf{Prinzip}: Speculative + Cache bleibt + Cache-Timing.


\begin{enumerate}
  \item Speculative greift auf Kernel-Adresse zu
  \item Daten in Cache
  \item Seg-Fault — Cache befüllt
  \item Flush\&Reload (Thread 2) misst Latenz
\end{enumerate}

\textbf{Probe-Array}:

\begin{lstlisting}
ProbeArray = [ASCII];
secret = readCharacter[5000];
character = ProbeArray[secret];
\end{lstlisting}

Schnellster Index = Geheimnis → iterieren.



\subsection{Reparatur}


{\small
\begin{adjustbox}{max width=\linewidth}
{\scriptsize
\begin{tabular}{ll}
\toprule
\textbf{Maßnahme} & \textbf{Wirkung} \\
\midrule
HW-Abschaltung OoO/Branch/Spec & drastischer Performance-Verlust \\
OS-Patch & \textasciitilde{}5–30\% Einbuße \\
\textbf{KASLR} & zufällige Kernel-Adressen \\
\textbf{KAISER/KPTI} & 2 AR (User vs. Kernel); Kernel nicht mehr in User-PT \\
Hardsplit (Zukunft) & feste HW-Hälften \\
AMD & keine Spekulation auf geschützten Bereich \\
\bottomrule
\end{tabular}
}
\end{adjustbox}
}


\begin{itemize}
  \item KASLR-Schwäche: Spraying + Side-Channel
  \item KAISER verlangsamt Syscalls + IRQ
\end{itemize}

\begin{quote}
Spectre deutlich schwerer zu patchen.
\end{quote}


\subsection{Recap}
\begin{itemize}
  \item Designfehler, kein SW-Bug
  \item Löschen/verschlüsseln keine Daten
  \item Nicht remote — Code muss lokal laufen
  \item Hebelt Isolation Prozess↔Prozess + Prozess↔Kernel auf
  \item Meltdown: leichter ausnutzbar + leichter behebbar als Spectre
\end{itemize}


\section{Geschichte (Eckdaten)}


{\small
\begin{adjustbox}{max width=\linewidth}
{\scriptsize
\begin{tabular}{ll}
\toprule
\textbf{Jahr} & \textbf{Ereignis} \\
\midrule
70 v.Chr. & Antikythera \\
1623 & Schickard \\
1642 & Pascaline (6-stellig Add.) \\
1801 & Jacquard-Webstuhl \\
1832 & Babbage; Lovelace (Bernoulli) \\
1887 & Hollerith → IBM \\
1905 & Vakuumdiode (Fleming) \\
1938/41 & Zuse Z1/Z3 \\
1944 & Harvard Mark I \\
1946 & ENIAC \\
1947 & Williams-Kilburn; Transistor \\
1958/59 & IC (Kilby/Noyce) \\
1964 & IBM OS/360 \\
1968 & Unix; Maus (Engelbart) \\
1969 & ARPANET; Apollo 11 \\
1971 & Intel 4004 \\
1976 & Apple I; IP \\
1981 & IBM 5150 \\
1985 & Windows 1.0 \\
1987/90 & Minix → Linux \\
1995 & Win95 (32-Bit) \\
2007 & iPhone \\
2023 & Frontier 1,194 EFLOPS \\
\bottomrule
\end{tabular}
}
\end{adjustbox}
}


\textbf{Fehler-Fälle}: Ariane 5 (64→16-Bit Overflow), Mars Climate Orbiter (SI vs. imperial Faktor 4,45), F-16 (negative Breitengrade), Spectre/Meltdown.







% ──────────────────────────────────────────────────────────────

\end{multicols}
\end{document}
